% Preparación separada del 30-09-2026; no modifica manuscritos publicados.
% Propietarios: VI_ES/sections/04_atlas_familias.tex:140--163;
% md/04d_ocupacion_estadistica.tex:305--319;
% md/06e_color_qutrit_gauge_finito.tex:1--194.
% Carta qutrit y forma compacta recibidas con sus hipótesis explícitas.
% Inserción: después de 00_antecedentes_geometricos_materiales y antes de 10.
\section{Residuo cromático y representación interna de la acción conjunta}
\label{hmtcuatro:color:seccion}

La cadena APP--TRIT--TPK conserva el estado enriquecido, sus rutas,
orientaciones, hojas y memoria dentro de la estructura discreta conjunta
del continuo. El residuo que se lee a continuación es una publicación
finita de ese estado; no lo sustituye ni reinicia su historia. Las
coordenadas de fase, incluida la unidad compleja y la coordenada regional
\(\pi\), son salidas HMT recibidas antes de esta realización algebraica.
La representación se construye sobre el soporte residual con la carta,
la orientación y el carácter de fase declarados. No se infiere un grupo
de calibre a partir del solo número de etiquetas.

\subsection{Soporte residual y orientación conservada}
\label{hmtcuatro:color:residuo}

En la copia posicional de APP, sean
\[
 I_9=\{1,\ldots,9\},\qquad
 \mathcal C=I_9\times I_9,\qquad
 \mathcal C_q=\{(r,c)\in\mathcal C:r\text{ es par}\}.
\]
El observable residual y la inversión celular son
\[
 \kappa_{\rm col}(r,c)=r-c\pmod3,\qquad
 \rho_{\rm c}(r,c)=(10-r,10-c).
\]
El símbolo \(\kappa_{\rm col}\) no designa el coeficiente gravitatorio
de la acción. La elevación de este observable a las secciones
persistentes se efectúa por la proyección al atlas y conserva los
restantes registros; no selecciona una celda por su masa evaluada.

\begin{proposition}[Fibras ternarias y balance cromático]
\label{hmtcuatro:color:balance}
El funcional \(\kappa_3:\mathbb F_3^2\to\mathbb F_3\),
\(\kappa_3(r,c)=r-c\), tiene tres fibras afines de cardinal tres.
Sobre \(\mathcal C_q\), los tres residuos tienen cardinales
\(12,12,12\). La inversión celular niega el residuo.
\end{proposition}
\begin{proof}
Para cada \(j\in\mathbb F_3\), la fibra es
\(\{(c+j,c):c\in\mathbb F_3\}\), una traslación del núcleo
\(\{(c,c)\}\). En \(\mathcal C_q\) existen cuatro valores pares
de \(r\); para cada uno, \(c=1,\ldots,9\) recorre tres veces
cada clase módulo tres. Cada residuo posee \(4\cdot3=12\)
representantes. Finalmente,
\[
 \kappa_{\rm col}(\rho_{\rm c}(r,c))
 =(10-r)-(10-c)=-(r-c)\pmod3.
\]
La clase cero queda fija y las clases uno y dos se intercambian.
\end{proof}

Este balance no identifica sabor, profundidad y color. Tampoco afirma
que las tres rectas que se construirán sean tres estados completos:
las rutas y la memoria que comparten una lectura residual permanecen
en la fibra enriquecida.

\subsection{Carta qutrit, par de Weyl y álgebra de color}
\label{hmtcuatro:color:weyl}

Se fija la carta
\[
 V_{\rm col}=\mathbb C[\mathbb F_3]\simeq\mathbb C^3,
 \qquad (e_0,e_1,e_2),\qquad
 \omega_3=e^{2\pi i/3},
\]
con base delta ortonormal, orientación \(j\mapsto j+1\) y carácter
\(j\mapsto\omega_3^j\). Los operadores son
\[
 Xe_j=e_{j+1\bmod3},\qquad Ze_j=\omega_3^j e_j,
 \qquad X^3=Z^3=I,\quad ZX=\omega_3XZ.
\]
La última igualdad se comprueba aplicando ambos miembros a cada
\(e_j\). La aplicación del atlas a esta carta es
\[
 (r,c)\longmapsto\mathbb C e_{\kappa_{\rm col}(r,c)}.
\]
Así, \(X\) permuta las rectas y \(Z\) registra sus fases. Cambiar
la carta afín y transportar conjuntamente base, orientación y carácter
conjuga los operadores por la correspondiente matriz unitaria de
cambio de base. No equivale a cambiar sólo una etiqueta y dejar fijos
los demás datos.

\begin{theorem}[Álgebra matricial y forma real compacta]
\label{hmtcuatro:color:algebra}
Los nueve operadores \(W_{ab}=X^aZ^b\),
\((a,b)\in\mathbb F_3^2\), forman una base ortogonal de
\(\operatorname{End}_{\mathbb C}(V_{\rm col})\). Los ocho índices
no nulos generan linealmente \(\mathfrak{sl}_3(\mathbb C)\), cuya
forma real de matrices antihermíticas es \(\mathfrak{su}(3)\).
\end{theorem}
\begin{proof}
Si \(a\ne0\), la matriz \(X^aZ^b\) no tiene entradas diagonales.
Si \(a=0\), su traza es \(1+\omega_3^b+\omega_3^{2b}\), que
vale cero para \(b\ne0\) y tres para \(b=0\). La relación de
Weyl da entonces
\[
 \operatorname{tr}(W_{ab}^{\dagger}W_{cd})
 =3\delta_{ac}\delta_{bd}.
\]
La independencia y la dimensión nueve prueban la primera afirmación;
los ocho restantes forman la base del subespacio de traza cero.
Además,
\[
 W_{ab}^{\dagger}=\omega_3^{ab}W_{-a,-b}.
\]
Los ocho índices no nulos forman cuatro pares adjuntos. Tomando
\(\mathcal R=\{(1,0),(0,1),(1,1),(1,2)\}\), las matrices
\[
 A_u=W_u-W_u^{\dagger},\qquad
 B_u=i(W_u+W_u^{\dagger}),\qquad u\in\mathcal R,
\]
son antihermíticas y de traza cero. Su independencia real se deduce
de la independencia compleja de los ocho \(W_u,W_{-u}\): en cada
par, los coeficientes de una combinación real nula fuerzan a cero
los dos coeficientes reales. Forman una base real de dimensión ocho.
Por tanto,
\[
 \{A\in\mathfrak{sl}_3(\mathbb C):-A^{\dagger}=A\}
 =\mathfrak{su}(3).
\]
\end{proof}

La forma hermítica y la orientación compleja fijadas determinan el
grupo de realización
\[
 SU(V_{\rm col})=\{U:U^{\dagger}U=I,\ \det U=1\}.
\]
Su álgebra tangente es exactamente la anterior: diferenciar ambas
ecuaciones en la identidad da \(A^{\dagger}=-A\) y
\(\operatorname{tr}A=0\); recíprocamente, \(e^{tA}\) satisface
ambas. Todo elemento del grupo se escribe como exponencial de una
matriz de esa álgebra. En efecto, diagonalícese unitariamente \(U\)
con autovalores \(e^{i\theta_j}\); como el determinante es uno,
\(\sum_j\theta_j=2\pi n\). Restar \(2\pi n\) a uno de los
ángulos no altera \(U\) y produce un logaritmo antihermítico de
traza cero. Esta realización continua no se identifica con el grupo
finito generado únicamente por \(X,Z\). Conserva la forma compacta
y no introduce un valor para el acoplamiento fuerte.

\subsection{Transporte cromático y factores internos independientes}
\label{hmtcuatro:color:transporte}

En una celularización finita orientada compatible con las rutas,
se asignan \(U_e\in SU(V_{\rm col})\) y
\(U_{\bar e}=U_e^{-1}\). Para \(g_v\in SU(V_{\rm col})\),
\[
 U_e^g=g_{t(e)}U_eg_{s(e)}^{-1},\qquad
 \psi_v^g=g_v\psi_v.
\]
Por sustitución,
\[
 U_e^g\psi_{s(e)}^g-\psi_{t(e)}^g
 =g_{t(e)}(U_e\psi_{s(e)}-\psi_{t(e)}).
\]
El signo opuesto, utilizado en \(D_T\), tiene la misma covariancia.
La transformación de un producto orientado alrededor de una cara
es conjugación por el calibre del vértice base: los factores de
los vértices intermedios se cancelan consecutivamente. En particular,
para \(\beta\ge0\),
\[
 S_{\rm col}^{(\mathcal K,\beta)}[U]
 =\frac\beta3\sum_f(3-\operatorname{Re}\operatorname{tr}U_f)
\]
es invariante y no negativa. La invariancia usa la traza cíclica;
la positividad resulta de \(\operatorname{Re}\operatorname{tr}U_f\le3\).
El parámetro y la celularización conservan su condición de datos de
esta realización. Esta prueba no los determina por comparación con
valores cromodinámicos.

La materia de la acción conjunta utiliza este factor, no otra
partición del atlas. Para los quarks izquierdos, su fibra interna
tiene factores \(V_{\rm col}\otimes\mathbb C^2\otimes F\),
donde \(F\) conserva las familias; para cada tipo derecho se usa
\(V_{\rm col}\otimes F\). El factor de color es trivial en los
leptones. El factor espinorial se conserva aparte. Para matrices
\(A\) de color y \(B\) débiles se cumple
\[
 [A\otimes I,I\otimes B]=0,
\]
pues los dos productos son \(A\otimes B\). Los caracteres de
hipercarga actúan como escalares sobre cada multiplete y también
conmutan. Esta identidad permite componer las conexiones internas
y la de espín en la misma derivada sin confundir color, quiralidad
y sabor.

Los mapas Yukawa de la acción posterior son la identidad en el
factor cromático y mapas en los factores Higgs--familia. Por ello
\(Y_q(H)(g_{\rm col}\otimes I)
=(g_{\rm col}\otimes I)Y_q(H)\). La compatibilidad débil y de
hipercarga se demuestra allí con sus columnas y sus pesos explícitos;
no se obtiene del censo \(12+12+12\).
Así se conserva la cadena residual--qutrit--álgebra--transporte
que requiere la acción común, sin afirmar en este antecedente un
límite cuántico de Yang--Mills ni reemplazar sus propietarios.
